\documentclass[10pt]{ctexart}
\usepackage[a4paper,margin=1in]{geometry}
\usepackage{amsmath,amssymb,mathtools,bm}
\usepackage{booktabs,longtable,array}
\usepackage{enumitem}
\usepackage{listings}
\usepackage{xcolor}
\usepackage{microtype}
\usepackage{hyperref}
\setlist[itemize]{leftmargin=1.5em}
\setlist[enumerate]{leftmargin=1.8em}
\lstset{basicstyle=\ttfamily\small,breaklines=true,columns=fullflexible,keepspaces=true,frame=single,framerule=0.2pt,xleftmargin=0.5em,xrightmargin=0.5em}
\newcommand{\clip}{\operatorname{clip}}
\newcommand{\uid}{\operatorname{uid}}
\newcommand{\trajid}{\operatorname{traj\_uid}}
\newcommand{\stepidx}{\operatorname{step\_idx}}
\title{015: StepPPO-v3 算法改动说明}
\author{}
\date{2026-06-09}
\begin{document}
\maketitle
\vspace{-1.2em}

本文说明从 StepPPO-v2 到 StepPPO-v3 需要做的算法和工程改动。目标是先实现一个清晰、可诊断、低干扰的 v3-core，而不是一次性加入所有可能 trick。

\section*{1. v2 的问题边界}

StepPPO-v2 的当前配置：

\begin{lstlisting}
algorithm.step_ppo.step_advantage_w: 0.5
algorithm.step_ppo.episode_mode: mean_norm
algorithm.step_ppo.final_advantage_mode: direct
algorithm.step_ppo.normalize_episode_advantage: false
algorithm.step_ppo.normalize_step_advantage: true
algorithm.step_ppo.normalize_final_advantage: false
actor_rollout_ref.actor.value_head.detach_value_backbone: false
actor_rollout_ref.actor.value_head.value_loss_coef: 0.03
actor_rollout_ref.actor.ppo_micro_batch_size_per_gpu: 8
\end{lstlisting}

已有日志暴露四个问题：

\begin{enumerate}
\item \texttt{detach\_value\_backbone=false} 时，value loss 会回传 actor backbone；GiGPO value-aux 已经证明仅增加 auxiliary value regression 就可能破坏原 GiGPO。
\item StepPPO-v2 的 episode advantage 可能读取 raw \texttt{episode\_rewards}，而 step reward 包含 invalid action penalty，导致 reward 口径不一致。
\item final advantage 同时混合 \texttt{A\_epi} 和 \texttt{A\_step}，但两者的尺度、归一化和 reward 语义不同。
\item policy loss 仍然是 token-level clipping，而 WebShop 的 environment action 是完整 response/action string。
\end{enumerate}

v3-core 对应地做四个改动：

\begin{lstlisting}
shared value backbone update  -> value-head-only update
raw episode reward mix        -> shaped step reward only
episode + step hybrid adv     -> pure step GAE advantage
token-level PPO ratio         -> action/sequence-level PPO ratio
\end{lstlisting}

\section*{2. 算法差异表}

\begin{small}
\begin{longtable}{@{}p{0.25\linewidth}p{0.32\linewidth}p{0.32\linewidth}@{}}
\toprule
项目 & StepPPO-v2 & StepPPO-v3-core \\
\midrule
MDP 粒度 & step advantage + token PPO loss & environment step / response action \\
reward source & \texttt{episode\_rewards} + \texttt{step\_rewards} & shaped step reward only \\
invalid penalty & step term 可见，episode term 可能不可见 & policy/value 全部可见 \\
episode advantage & GiGPO-style \texttt{mean\_norm} & 删除 \\
step advantage & GAE + batch whitening & GAE + uid group whitening \\
final advantage & \texttt{A\_epi + 0.5 * A\_step} & \texttt{A\_step\_group\_norm} \\
value head & shared head, no detach & shared head, detach backbone \\
value loss & token-shaped/broadcast path & step scalar value loss \\
policy ratio & token-level PPO ratio & sequence geometric ratio \\
loss aggregation & token-oriented & \texttt{seq-mean-token-mean} \\
primary risk & value/backbone interference, reward mismatch & value underfit, lower update strength \\
\bottomrule
\end{longtable}
\end{small}


\section*{3. New Config Namespace}

建议新增独立 namespace，避免继续复用 v2 的 \texttt{algorithm.step\_ppo} 语义：

\begin{lstlisting}
algorithm:
  adv_estimator: step_ppo_v3
  gamma: 0.95
  step_ppo_v3:
    reward_source: shaped_step
    use_episode_advantage: false
    step_gamma: 0.95
    step_lam: 0.90
    advantage_norm: uid
    zero_no_variance_group: true
    policy_ratio: sequence_geometric
    final_advantage_mode: pure_step
    min_group_std: 1.0e-6

actor_rollout_ref:
  actor:
    policy_loss: action_level
    loss_agg_mode: seq-mean-token-mean
    value_head:
      enable: true
      value_position: pre_action_last_context_token
      detach_value_backbone: true
      value_loss_coef: 0.03
      cliprange_value: 0.5
\end{lstlisting}

保留 \texttt{algorithm.gamma=0.95} 是为了和当前 GiGPO WebShop 口径对齐。\texttt{step\_lam=0.90} 是 v3 的新默认值，用于降低 step return 方差。

\section*{4. Advantage 计算改动}

\subsection*{4.1 删除 episode score 输入}

v2 逻辑中：

\begin{lstlisting}
if "episode_rewards" in data.non_tensor_batch:
    episode_scores = data.non_tensor_batch["episode_rewards"]
else:
    episode_scores = token_level_rewards.sum(dim=-1)
\end{lstlisting}

v3 不再把 \texttt{episode\_scores} 传入 policy advantage。\texttt{episode\_rewards} 可以继续作为 logging 字段，但不能进入 \texttt{A\_v3}。

\subsection*{4.2 统一 shaped step reward}

新增 helper：

\begin{lstlisting}
def build_shaped_step_rewards(data, config) -> torch.Tensor:
    rewards = raw_rewards
    rewards -= invalid_action_penalty_coef * (1 - is_action_valid)
    if use_kl_in_reward:
        rewards += (token_level_rewards - token_level_scores).sum(dim=-1)
    return rewards
\end{lstlisting}

该 helper 应同时服务：

\begin{lstlisting}
step GAE advantage
step return target
reward diagnostics
\end{lstlisting}

不要再出现 policy advantage 使用 raw reward、value target 使用 penalized reward 的分裂。

\subsection*{4.3 新增 v3 GAE 输出}

建议新增函数：

\begin{lstlisting}
def compute_step_ppo_v3_advantage_return(
    step_rewards: torch.Tensor,
    state_values: torch.Tensor,
    uid: np.ndarray,
    traj_index: np.ndarray,
    step_index: np.ndarray,
    gamma: float,
    lam: float,
    sample_index: np.ndarray | None,
    norm: str = "uid",
    min_group_std: float = 1e-6,
):
    ...
    return advantages, step_returns, diagnostics
\end{lstlisting}

输出：

\begin{lstlisting}
advantages: shape [num_step_samples], scalar per environment action
step_returns: shape [num_step_samples], scalar per environment action
diagnostics:
  raw_adv_mean/std
  norm_adv_mean/std
  no_variance_group_ratio
\end{lstlisting}

注意：\texttt{compute\_step\_gae\_advantage\_return} 当前已经处理了 \texttt{step\_sample\_uid} 去重逻辑，v3 应复用这部分能力，不要重新引入 duplicated row 被当作 fake transition 的问题。

\subsection*{4.4 uid group whitening}

新增 normalization：

\begin{lstlisting}
for each uid group:
    group_adv = raw_adv[group_indices]
    if group_adv.numel() <= 1 or group_adv.std() < min_group_std:
        normalized[group_indices] = 0
        no_variance_group_count += 1
    else:
        normalized[group_indices] = (group_adv - group_adv.mean()) / (group_adv.std() + eps)
\end{lstlisting}

理由：

\begin{itemize}
\item 避免不同 task difficulty 混在一个 batch 里互相缩放。
\item 保留 GiGPO/GRPO 的 group-relative 稳定性。
\item 对全成功或全失败、没有相对差异的 group 自动降权。
\end{itemize}

\section*{5. Policy Loss 改动}

\subsection*{5.1 v2 token-level PPO}

v2 当前 token-level ratio：

\begin{lstlisting}
ratio = exp(log_prob - old_log_prob)
pg_loss = min(ratio * A_token, clip(ratio) * A_token)
\end{lstlisting}

同一个 response 内每个 token 单独 ratio、单独 clipping。

\subsection*{5.2 v3 sequence geometric ratio}

v3-core 使用 response/action-level ratio：

\begin{lstlisting}
token_log_ratio = (log_prob - old_log_prob) * response_mask
seq_len = response_mask.sum(dim=-1).clamp(min=1)
seq_log_ratio = token_log_ratio.sum(dim=-1) / seq_len
seq_ratio = exp(seq_log_ratio)
\end{lstlisting}

然后：

\begin{lstlisting}
pg_loss_i = -min(
    seq_ratio_i * A_i,
    clip(seq_ratio_i, 1 - eps_low, 1 + eps_high) * A_i,
)
\end{lstlisting}

实现选择：

\begin{enumerate}
\item 最小改动：复用已有 \texttt{compute\_policy\_loss\_gspo}，并设置 \texttt{actor.policy\_loss=gspo} 或新增 alias \texttt{action\_level}。
\item 更清晰改动：新增 \texttt{compute\_policy\_loss\_action\_level}，返回 action-level metrics，不复用 GSPO 命名。
\end{enumerate}

建议先走第 2 条，避免 method 文档和代码里混用 GSPO 名称。

\subsection*{5.3 必加 action-level metrics}

新增：

\begin{lstlisting}
actor/action_log_ratio_mean
actor/action_log_ratio_std
actor/action_ratio_mean
actor/action_ratio_std
actor/action_clipfrac
actor/action_approx_kl
\end{lstlisting}

token-level metrics 可以保留，但不能只看 token-level \texttt{ppo\_kl}。v3 要判断的是完整 response/action 是否被 PPO trust region 约束住。

\section*{6. Value Loss 改动}

\subsection*{6.1 detach backbone 默认开启}

v3 script 默认：

\begin{lstlisting}
actor_rollout_ref.actor.value_head.detach_value_backbone: true
\end{lstlisting}

代码已经支持该字段，\texttt{value\_head.py} 中会在 value head 前 detach hidden states。因此 v3 不需要先改 value head forward。

\subsection*{6.2 scalar value loss}

当前 \texttt{step\_returns} 在实现里容易沿 response token 维度 broadcast。v3 建议改成 step scalar：

\begin{lstlisting}
current_values: [batch]
old_values: [batch]
step_returns: [batch]
\end{lstlisting}

loss：

\begin{lstlisting}
values_clipped = old_values + clamp(current_values - old_values, -cliprange_value, cliprange_value)
value_loss_unclipped = (current_values - step_returns).pow(2)
value_loss_clipped = (values_clipped - step_returns).pow(2)
value_loss = 0.5 * max(value_loss_unclipped, value_loss_clipped).mean()
\end{lstlisting}

这和现有 \texttt{StepWisePPOActor} 的核心公式一致，但 v3 要确保 value loss 是 per action/step，而不是按 response token 重复加权。

\subsection*{6.3 value diagnostics}

新增：

\begin{lstlisting}
value_error = current_values.detach() - step_returns
value_rmse = sqrt(mean(value_error ** 2))
value_mae = mean(abs(value_error))
explained_variance = 1 - var(step_returns - pred) / (var(step_returns) + eps)
\end{lstlisting}

并记录分布：

\begin{lstlisting}
actor/value_pred_std
actor/value_return_std
actor/value_return_p05/p50/p95
\end{lstlisting}

invalid action 分层：

\begin{lstlisting}
actor/value_loss_valid
actor/value_loss_invalid
critic/adv_valid_mean
critic/adv_invalid_mean
\end{lstlisting}

\section*{7. DataProto 字段约定}

v3 建议保持字段清晰：

\begin{lstlisting}
data.batch["advantages"]      token-shaped or step-shaped policy advantages
data.batch["step_advantages"] scalar [batch] final v3 step advantage
data.batch["step_returns"]    scalar [batch] value target
data.batch["value_step_rewards"] scalar [batch] shaped step rewards
data.batch["state_values"]    scalar [batch] old state values
\end{lstlisting}

如果 actor update path 仍要求 token-shaped \texttt{advantages}，可以在最后一步 broadcast：

\begin{lstlisting}
advantages_token = step_advantages.unsqueeze(-1) * response_mask
\end{lstlisting}

但 \texttt{step\_advantages} 必须保留为 scalar 字段，用于 action-level loss 和 diagnostics。

\section*{8. 需要修改的文件}

\subsection*{8.1 \texttt{verl/trainer/ppo/step\_ppo\_algos.py}}

新增：

\begin{lstlisting}
build_shaped_step_rewards
normalize_advantage_by_uid
compute_step_ppo_v3_advantage_return
compute_step_ppo_v3_advantage_data
\end{lstlisting}

或在现有 \texttt{compute\_step\_ppo\_advantage\_data} 中通过 \texttt{algorithm.adv\_estimator == "step\_ppo\_v3"} dispatch。

\subsection*{8.2 \texttt{verl/trainer/ppo/core\_algos.py}}

新增：

\begin{lstlisting}
compute_policy_loss_action_level
\end{lstlisting}

该函数可以从 \texttt{compute\_policy\_loss\_gspo} 拆出来，但需要返回 action-level metrics。

\subsection*{8.3 \texttt{verl/workers/actor/step\_ppo\_actor.py}}

改动：

\begin{lstlisting}
支持 actor.policy_loss == "action_level"
读取 scalar step_advantages / step_returns
计算 action-level loss
记录 action-level ratio metrics
记录 value RMSE/MAE/EV/quantile metrics
\end{lstlisting}

\subsection*{8.4 \texttt{verl/trainer/config/step\_ppo\_trainer.yaml}}

新增：

\begin{lstlisting}
algorithm:
  step_ppo_v3:
    reward_source: shaped_step
    use_episode_advantage: false
    step_gamma: 0.95
    step_lam: 0.90
    advantage_norm: uid
    zero_no_variance_group: true
    min_group_std: 1.0e-6
    policy_ratio: sequence_geometric
    final_advantage_mode: pure_step
\end{lstlisting}

actor config 增加：

\begin{lstlisting}
policy_loss: token_level  # choices: token_level, gspo, action_level
\end{lstlisting}

\subsection*{8.5 \texttt{EXPS/run\_webshop\_step\_ppo\_v3\_4gpu\_paper\_align\_simple.sh}}

新增主实验脚本，建议基于 v2 脚本改：

\begin{lstlisting}
bash "$BASE_SCRIPT" vllm \
  algorithm.adv_estimator=step_ppo_v3 \
  algorithm.gamma=0.95 \
  algorithm.step_ppo_v3.step_gamma=0.95 \
  algorithm.step_ppo_v3.step_lam=0.90 \
  algorithm.step_ppo_v3.advantage_norm=uid \
  algorithm.step_ppo_v3.use_episode_advantage=false \
  actor_rollout_ref.actor.policy_loss=action_level \
  actor_rollout_ref.actor.loss_agg_mode=seq-mean-token-mean \
  actor_rollout_ref.actor.value_head.detach_value_backbone=true \
  actor_rollout_ref.actor.value_head.value_loss_coef=0.03 \
  actor_rollout_ref.actor.ppo_micro_batch_size_per_gpu=8 \
  trainer.experiment_name=step_ppo_v3_qwen2.5_1.5b_4gpu_paper_align_full_e250 \
  "$@"
\end{lstlisting}

\section*{9. Smoke Test Checklist}

2GPU smoke 必须先检查：

\begin{lstlisting}
1. 日志出现 actor/action_ratio_mean 和 actor/action_clipfrac。
2. actor/value_rmse、actor/value_explained_variance 非 NaN。
3. actor/value_loss 有值，但 detach=true 下 actor grad norm 不因 value loss 明显放大。
4. episode/valid_action_ratio 不应长期低于 0.5。
5. response_length/clip_ratio 不应快速升高到 v2 后期水平。
6. no_variance_group_ratio 合理；如果长期接近 1，说明 uid group norm 没有有效 policy signal。
\end{lstlisting}

如果 smoke 中出现 \texttt{action\_ratio\_mean} 爆炸：

\begin{lstlisting}
确认用的是 sequence geometric ratio，而不是 raw exp(sum log-ratio)。
确认 response_mask 没有包含 padding。
确认 loss_agg_mode 是 seq-mean-token-mean。
\end{lstlisting}

如果 value metrics 全部正常但 valid action ratio 仍低：

\begin{lstlisting}
优先检查 reward parser / invalid penalty 是否进入 shaped step reward。
再跑 v3_no_action_ratio，确认 action-level ratio 是否改变了 behavior。
\end{lstlisting}

\section*{10. Ablation Matrix}

首轮不要铺太多实验，建议：

\begin{small}
\begin{longtable}{@{}p{0.25\linewidth}p{0.32\linewidth}p{0.32\linewidth}@{}}
\toprule
实验 & 改动 & 目的 \\
\midrule
v3\_core & 默认 v3 & 主线 \\
v3\_no\_action\_ratio & \texttt{policy\_loss=token\_level} & 判断 action-level clipping 是否必要 \\
v3\_detach\_false & \texttt{detach\_value\_backbone=false} & 复验 shared-backbone interference \\
v3\_lam\_1 & \texttt{step\_lam=1.0} & 判断高方差是否来自 MC target \\
v3\_batch\_norm\_adv & \texttt{advantage\_norm=batch} & 判断 uid norm 是否必要 \\
\bottomrule
\end{longtable}
\end{small}


对照必须包含：

\begin{lstlisting}
GiGPO seed 2026
StepPPO-v2 seed 2026
GiGPO value-aux detach=true seed 2026
\end{lstlisting}

\section*{11. 成功/失败判据}

\subsection*{11.1 稳定性判据}

v3\_core 最低稳定性标准：

\begin{lstlisting}
valid_action_ratio >= 0.85 after warmup
response_clip_ratio close to GiGPO order of magnitude
action_clipfrac not persistently > 0.3
value_loss no repeated spikes above 8
no validation collapse similar to value-aux step 180--185
\end{lstlisting}

\subsection*{11.2 效果判据}

如果 v3\_core step 200 仍显著低于 StepPPO-v2，先不要继续加复杂 critic，优先定位：

\begin{lstlisting}
reward口径是否一致
step_lam是否过低导致 credit 传播不足
uid whitening是否把全部正向信号消掉
action-level ratio是否过度保守
\end{lstlisting}

如果 v3\_core 稳定但只接近 StepPPO-v2：

\begin{lstlisting}
尝试 step_lam=0.95
尝试 value_loss_coef=0.01/0.05
尝试 separate value head lr
尝试 GAGPO-style grouped value proxy 替代 learned value target
\end{lstlisting}

\section*{12. 不纳入 v3-core 的内容}

以下内容不进入 v3-core，避免主线过宽：

\begin{enumerate}
\item separate critic model。
\item value loss 回传 backbone。
\item in-distribution critic pretraining。
\item Huber value loss 默认开启。
\item PRM / external reward model。
\item GAGPO grouped value proxy 替代 learned value。
\item dynamic task sampling。
\end{enumerate}

这些方向都可能有价值，但应在 v3-core 稳定后作为 v3.1/v4 单独验证。

\section*{13. Implementation Order}

建议实现顺序：

\begin{enumerate}
\item 新增 v3 advantage 计算和 shaped reward helper。
\item 新增 action-level policy loss 和 metrics。
\item 改 StepWisePPOActor 支持 scalar \texttt{step\_advantages} / \texttt{step\_returns}。
\item 增加 value diagnostics。
\item 增加 config namespace 和 v3 run script。
\item 跑 2GPU smoke。
\item 跑 seed 2026 full 250 step。
\item 对比 GiGPO、StepPPO-v2、GiGPO value-aux detach=true。
\end{enumerate}

每一步完成后都应保留可回退的独立开关，避免 v3 出问题时只能整体回滚。

\end{document}
